\subsection{Structure at the neuron level does not imply structure at another scale}

  % \caption{\textbf{A geometric framework for comparing neural populations at scale.}
  % \textbf{a.} Shape metrics compare representations across animals performing the same task,
  % across regions within one animal, or across sessions of one animal over time.
  % \textbf{b.} Population activity over $M$ conditions and $N$ neurons, here a periodic
  % variable $\theta$ (e.g.\ head direction). Tuning curves (\textit{left}) are $M$ landmarks in
  % $N$-dimensional firing-rate space (\textit{right}), tracing a manifold; colour denotes
  % $\theta$.
  % \textbf{e.} The distance $d(\mbX_i, \mbX_j)$ is the dissimilarity between two manifolds
  % after alignment.
  % \textbf{f.} Manifolds and this distance form a \textit{metric space}, in which clustering
  % and regression carry their usual guarantees.
  % \textbf{c--h.} Simulated fields of $K = 60$ regions, $N = 100$ neurons each. Every neuron is
  % a point in a two-dimensional space of tuning shapes (skew, width) and a region is a cloud in
  % that space: $a$ sets how bimodal the cloud is along its long axis, $\psi$ the orientation of
  % that axis. The two are independent by construction, and neither changes a region's mean
  % tuning or total spread.
  % \textbf{c.} Example regions varying in $a$ alone (\textit{top}) or $\psi$ alone
  % (\textit{bottom}); ellipses are 2 s.d.\ Gaussian fits to each cluster.
  % \textbf{d.} Neuron-level analyses of two fields: one whose regions vary along a gradient
  % (\textit{top}) and one where the same regions are arranged at random (\textit{bottom}).
  % \textit{Left}, per-region categoricality, silhouette against a matched Gaussian, shading
  % $|z| < 1.96$; \textit{middle}, the same test on neurons pooled across regions
  % ($z = 6.0$ and $6.4$); \textit{right}, regions embedded by their averaged tuning curve. The
  % two fields are indistinguishable throughout.
  % \textbf{g.} The same fields, plus one built from three discrete region types, embedded by
  % classical MDS of their $K \times K$ Procrustes distances; colour denotes each region's
  % categoricality.
  % \textbf{h.} Best silhouette between regions against Gaussians matched to each embedding
  % ($p = 0.199$, $0.557$, $0.005$).}
  
\begin{figure}[p]
    \centering
    \includegraphics[width=1\textwidth]{content/figures/netrep-journal-figure-1.pdf}

       \caption{\textbf{A geometric framework for comparing neural populations across scales.}
    \textbf{a.} Shape metrics can be leveraged to compare representations across animals, regions or different time periods from the same animal. \textbf{b.} Example population activity measured across $M$ conditions and $N$ neurons in a simulated experiment parameterized by a periodic 1D variable, $\theta$ (e.g. head direction). Smooth tuning curves, \textit{left}, can be viewed as a smooth manifold, \textit{right}, in $N$-dimensional firing rate space. Color code denotes values of $\theta$. \textbf{c} Simulation of 3 regions, $N = 100$ neurons each where population geometry and clusters in neuron space are decoupled. Every neuron is a point in a two-dimensional space of tuning shapes (skew, width), shown here as the first two principal components of the neurons $\times$ conditions matrix, with the skew and width directions recovered by regression (grey axes). \textit{Left}, a region without any clusters. \textit{Middle}, another region without any clusters but a rotated cloud in width-skew space leads to a different geometry, bottom. \textit{Right}, a region with the same overall orientation as on the left, but with two clusters, leads to a similar geometry. Bottom row, each region's population manifold, with the left region's manifold repeated in grey behind all three after optimal alignment. \textbf{d.} Neuron-level analyses of two population of regions that do not cluster in region space: one whose regional geometry varies along a gradient coupled with strength of within-region clusters (\textit{top}) and one where geometries are decoupled from within-region clusters (\textit{bottom}). \textit{Left}, per-region categoricality (\cite{posani2026}), silhouette against a matched Gaussian, shading $|z| < 1.96$; \textit{right}, the same test on neurons pooled across regions ($z = 6.0$ and $6.4$); In both cases, regions are inferred to be clustered. \textbf{e.} The goal of our framework isto compute a distance function $d(X_i,X_j)$, which we define as the dissimilarity between two manifolds $X_i, X_j$ after alignment. \textbf{f.} The set of possible manifolds and the distance function together define a \textit{metric space}. The properties of these metric spaces allow us to leverage existing methods, such as clustering and regression with rigorous theoretical guarantees. \textbf{g.} The same population of regions, plus one built from three discrete region types, embedded by MDS of their $K \times K$ Procrustes distances; colour denotes each region's categoricality. Below each, the best silhouette between regions against a Gaussian matched to the embedding: only the field built from discrete types exceeds its null. \textbf{h.} Both analyses over a two-dimensional design in which within-region clustering (vertical) and across-region clustering (horizontal) are varied independently; colour is $z$ against each method's own null, on independent scales  per panel. Pooling neurons responds to within-region clustering and is flat along the other axis; the region-space test does the reverse.
    }
        \label{fig:schematic}
\end{figure}

% Before turning to experimental data, we show with simulated data how challenging it is to unveil structure across systems and how a proper metric can be leveraged to unveil such structure. 

Consider a set of neural recordings collected across $K$ biological systems---for example, $K$ different animals, $K$ behavioral sessions from the same animal, or recordings from $K$ different brain regions (\textbf{Fig.}~\ref{fig:schematic}a).
The activity of each neural system is recorded at $M$ measurement points, corresponding to different experimental conditions or time points.
For example, neural responses could be measured in response to $C$ conditions (e.g. distinct sensory stimuli or trained motor actions) over $T$ time points per condition.
This would result in a total of $M = C\,T$ measurements.
To simplify the exposition, we will assume that each neural population is composed of $N$ neurons.\footnote{This assumption is typically violated in practice, but the approach is easy to adapt to unequal neural populations (see \textit{Methods} and \cite{harvey2024duality}).} In this setting, the data collected from $K$ neural systems are represented by a set of $N \times M$ matrices, $\mbX_1, \mbX_2, \dots, \mbX_K$.
Each row corresponds to a neuron's tuning curve across the $M$ measurement points (task conditions and time points).
Each column corresponds to an $N$-dimensional population response to one of the experimental measurements.
\textbf{Fig.}~\ref{fig:schematic}b illustrates an example where the measurements span a 1-dimensional space, resulting in a 1-dimensional, smooth manifold structure within the neural firing rate space.

For the sake of concreteness, we will now focus on how to reveal structure across $K$ `regions' in simulation, but these principles apply directly to other scales (\textbf{Fig. } \ref{fig:schematic}a). We simulate different regions as a population of $N = 100$ neurons that collectively encode a periodic variable $\theta$ (\textbf{Fig. } \ref{fig:schematic}b). All tuning curves peaks at the same angle, but each neuron tuning curve varies in their width and skewness ($\kappa$ and $s$, respectively) of von Mises curves (Methods). A neuron can therefore be described by two numbers, how skewed and how broad its tuning is (\textbf{Fig.} \ref{fig:schematic}c). Within each simulated region, neurons' tuning was structured differed according to two parameters ($\psi$ and $a$). On one hand, $a$ sets the number of clusters in skew-width space: at $a = 0$ neurons form a single continuous cloud of tuning shapes (\textbf{Fig. }\ref{fig:schematic}c, \textit{left}) and as $a \to 1$ they separate into two discrete clouds (\textbf{Fig. }\ref{fig:schematic}c, \textit{right}). In other words, $a$ sets how categorical those neurons are \parencite{posani2026}. In turn, $\psi$ sets the orientation of these clouds.  At $\psi = 0$ neurons differ in $s$ while sharing a common $\kappa$, at $\psi = 90^\circ$ they differ in $\kappa$ while sharing a common skew $s$, and at intermediate angles $s$ and $\kappa$ co-vary (\textbf{Fig. }\ref{fig:schematic}c, \textit{dashed line}). Importantly, $a$ does not affect a region's population geometry, whereas $\psi$ does. This holds only because we fixed single-neuron selectivity (\textbf{Fig. } \ref{fig:schematic}c, \textit{bottom}).

% Along the long axis, tuning parameters were drawn from a mixture of two components centred at $\pm a \sigma$ with within-component spread $\sigma \sqrt{1 - a^2}$, so that the variance along that axis equals
% $\sigma^2$ for every value of $a$.
% At $a = 0$ a region is a single continuous cloud of tuning shapes; as $a \to 1$ it separates into two discrete cell types; and the total diversity of tuning is identical throughout.
% Crucially, this construction leaves every region with the same mean tuning curve and the same mean tuning strength.
% Regions differ only in how their tuning diversity is \textit{organised}, not in how much of it there is.

We then simulated three different scenarios regarding structure across regions. These scenarios were build to exemplify different extremes that are particularly hard to distinguish by clustering neurons.
%: one in which neurons within regions might be categorical, but regions themselves are not; and another in which neurons are not categorical at all, but regions are strongly categorical. For the former case, we simulate two different possibilities: one in which regions were organized in a gradient and another in which the were not. Specifically, we simulated three cohorts of $K = 60$ regions.
In the first, which we called \textit{gradient}, regions were sampled along a smooth gradient in which $a$ and $\psi$ varied together, so that some regions are strongly categorical and others are not. In the second, which we called \textit{unstructured}, we also varied $a$ and $\psi$ across regions but broke their paring. This also lead to very different structure within regions, but not organized in a gradient across them. Finally, in the third scenario no region is categorical at all ($a = 0$ throughout), but $\psi$ took one of three discrete values, giving three types of region. Note that in these scenarios structure within and across regions is decoupled: neurons within regions might be categorical, but regions themselves are not (\textit{graded} and \textit{unstructured)}; conversely, neurons might not categorical at all, but regions can be strongly categorical (\textit{clustered}). 

This setup simulates an important question in neuroscience: can we distinguish between a group of regions that falls into discrete \textit{clusters} from those that do not, either in a \textit{gradient} or \textit{unstructured}? One recently proposed strategy is to quantify the structure within regions using by clustering single-neuron response properties \parencite{posani2026}. To show how challenging it is to tackle the aforementioned question at the single-neuron level, we then quantify if neurons formed clusters within a region using the pipeline of \textcite{posani2026}: each neuron's tuning curves clusters are inferred with $k$-means and compared against a null model without clusters (see \textit{\nameref{sec:methods}} and \textcite{posani2026}). As can be seen in \textbf{Fig.} \ref{fig:schematic}d, \textit{left}, both gradient and unstructured scenarios had similar distributions of clusters ($z$). A previously proposed strategy to quantify structure \textit{across} regions is to cluster neurons pooling across all regions (\textcite{posani2026}). However, this tests if \textit{within} region structure is different across regions, but is blind to \textit{across} region structure. Indeed, applying this method to our simulations shows a significant deviation from the null model for both gradient (\textbf{Fig.} \ref{fig:schematic}d, \textit{top right}) and unstructured scenario (\textbf{Fig.} \ref{fig:schematic}d, \textit{bottom right}).

In this manuscript, we argue that open questions of this kind are readily answerable with existing methods, provided that regions --- or individuals, recording sessions, etc. --- are first placed in a common space at the appropriate scale, within which standard tools such as clustering and regression directly apply. 

\subsection{Comparing neural populations as points in a metric space}


% As a simple first example, we will consider tasks that produce smooth, periodic tuning curves encoded by a task parameter $\theta \in [-\pi, \pi)$, which is discretized into $M$ bins.
% This is a generic form of tuning that pervades the neuroscience literature.
% Examples include representations of heading direction (e.g. \cite{chaudhuri,duszkiewicz}), orientation of a visual grating (e.g. \cite{hubel}), target locations in radial reaching tasks (e.g. \cite{georgopoulos}), and saccade targets in working memory tasks (e.g \cite{funahashi}).
% In all of these cases, the neural firing rates will trace a smooth ring manifold in $N$-dimensional firing rate space (\textbf{Fig.}~\ref{fig:schematic}d).

Our first objective is to define a family of distance functions $d(\mbX_i, \mbX_j)$ that quantify the dissimilarity between two neural systems.
The principal challenge is that the neurons are randomly sampled and labeled in each system and so cannot be immediately compared.
This can be interpreted geometrically as two neural manifolds being randomly oriented with respect to each other in firing rate space (\textbf{Fig.}~\ref{fig:schematic}e).
Although neural dimensions are not matched between systems, the $M$ measurements are shared (coloring in \textbf{Fig.}~\ref{fig:schematic}d-e).
Thus, we can consider \textit{aligning} the two manifolds through some transformation and then computing some measure of distance (\textbf{Fig.}~\ref{fig:schematic}f). 
% This is the motivating intuition for the shape metrics framework as well as several existing methods (e.g. linear predictivity scores and CCA).
% Intuitively, difference choices for the alignment transformation lead to different notions of distance between systems.

We will argue that it is useful to focus on distances which satisfy the following three axiomatic properties:
\begin{align}
\label{eq:metric-space-equivalence}
\text{Equivalence:} & \quad d\left(\mbX_i, \mbX_j\right)=0 \qquad \text{if and only if $\mbX_i$ and $\mbX_j$ are ``perfectly alignable''}\\
\label{eq:metric-space-symmetry}
\text{Symmetry:} & \quad d\left(\mbX_i, \mbX_j \right) =d\left(\mbX_j, \mbX_i \right)\\
\label{eq:metric-space-triangle}
\text{Triangle Inequality:} & \quad d\left(\mbX_i, \mbX_j \right) \leq d\left(\boldsymbol{X}_i, \mbX_k \right) + d\left(\mbX_k, \mbX_j \right)
\end{align}
where $\mbX_i$ and $\mbX_j$ are said to be ``perfectly alignable'' if the transform class specified by the practitioner is flexible enough to superimpose the two manifolds.

Notions of distance that satisfy \cref{eq:metric-space-equivalence,eq:metric-space-symmetry,eq:metric-space-triangle} are said to define a \textit{metric space}.
Intuitively, metric spaces are useful because they enable us to spatially reason about data.
That is, we can envision each of our recordings from $K$ neural systems $\mbX_1, \dots, \mbX_K$ as points within the space of all possible neural recordings (\textbf{Fig.}~\ref{fig:schematic}f).
That is, in this (potentially high-dimensional) metric space, an entire neural system is represented by a single point, and the distances between points indicate the geometric similarity or dissimilarity of the systems.   
These properties allow us to leverage methods such as clustering and nearest-neighbor regression (\textbf{Fig.}~\ref{fig:schematic}g) and to do so with theoretical guarantees \parencite[see e.g.][]{Cover1967,Dasgupta2005}.

The \textit{Procrustes shape distance} \parencite{schonemann1966generalized} is a metric that satisfies \cref{eq:metric-space-equivalence,eq:metric-space-symmetry,eq:metric-space-triangle}, and most of our analyses will leverage this distance as a running example of the shape metrics framework.
If we assume that each neuron's tuning curve has been pre-processed to be mean centered (i.e. rows of $\mbX_i$ and $\mbX_j$ have mean zero), then the Procrustes distance can be written as:
\begin{equation}
\label{eq:procrustes-distance}
d_{\text{Proc}}(\mbX_i, \mbX_j) = \min_{\mbQ \in \cO(N)} \Vert \mbQ \mbX_i - \mbX_j \Vert_F
\end{equation}
Here, $\cO(N)$ is the set of $N \times N$ orthogonal matrices.
Intuitively, $\mbQ$ encodes the optimal rotation and reflection that minimizes the sum of squared residuals between the two neural manifolds.
The mean centering step ensures that the rotation and reflection are performed around the manifold's center of mass.

With a proper metric in hand, each of the $K$ recordings becomes a single point in shape space, and the distances between them become the object of study (\textbf{Fig.}~\ref{fig:schematic}f). This makes the question we posed above answerable. Returning to the three simulated scenarios, we computed the $K \times K$ matrix of Procrustes distances, $\mbD_{ij} = d_{\text{Proc}}(\mbX_i, \mbX_j)$, for each scenario allowing us to visualize the \textit{regions space} with multidimensional scaling (\textbf{Fig.}~\ref{fig:schematic}g, \textit{top}). This reveals clearly different structures across the three scenarios: the \textit{gradient} scenario (\textbf{Fig.}~\ref{fig:schematic}g, \textit{left}) forms a single cloud through which categoricality varies smoothly ($\rho = 0.72$), the \textit{unstructured} scenario (\textbf{Fig.}~\ref{fig:schematic}g, \textit{middle}) forms a similar same cloud but with categoricality scattered across it at random ($\rho = 0.06$), and the \textit{clustered} scenario (\textbf{Fig.}~\ref{fig:schematic}g, \textit{right}) forms three well-separated groups. We can test this formally directly in region space by comparing the silhouette between regions against Gaussians matched to each embedding. Only the third is lumpier than a single continuous cloud ($p = 0.005$; \textbf{Fig.}~\ref{fig:schematic}h).
% Establishing a metric space enables us to use a broad set of off-the-shelf analysis tools.
% These methods take in a matrix of pairwise distances $\mbD_{ij} = d(\mbX_i, \mbX_j)$ and output unsupervised or supervised summaries of variability across the neural systems.
% For example, we may be interested in studying individual variability across $K$ animal subjects.
% To identify clusters of subjects with similar neural tuning, we may use hierarchical clustering methods, which work with theoretical guarantees in metric spaces~\parencite{Dasgupta2005}.
% To visualize these clusters, we can use multidimensional scaling (MDS) and other vector embedding techniques~\parencite{Borg2005-gk}.
% Further, if we measure a behavioral variable for each animal (e.g. performance on a task), we can use the distance relations between animals, based on neural data, to fit models that predict this behavior.
% For example, nearest nearest neighbor regression works with theoretical guarantees in metric spaces~\parencite{Cover1967}, and we will later demonstrate that it works well in practice.

% How can we define meaningful notions of distance between neural systems?
% This problem is nontrivial because simple measures, such as the Euclidean distance:
% \begin{equation}
% \Vert \mbX_i - \mbX_j \Vert_F = \left ( \sum_{n=1}^{N} \sum_{m=1}^{M} \left(\mbX_i[n, m] - \mbX_j[n, m] \right)^2 \right )^{1/2}
% \end{equation}
% are clearly meaningless whenever neurons are not identified and one-to-one matched across neural systems (\hl{Figure}).
% Intuitively, we need some way of \textit{aligning} the neural representations across two systems.
% We propose to use a constrained set of linear transformations, parameterized by a matrix $\mbQ$, to perform this alignment.
% Further, since it is common to pre-process neural data (e.g. by z-scoring each neuron) we propose to apply a user-defined function $\phi(\cdot)$ before fitting the alignment.
% Putting these pieces together, we can define:
% \begin{equation}
% \label{eq:generalized-euclidean-shape-metric}
% d(\mbX_i, \mbX_j) = \min_{\mbQ \in \cG} \Vert \mbQ \phi(\mbX_i) - \phi(\mbX_j) \Vert_F
% \end{equation}
% where $\cG$ denotes the set of allowable linear transformations.


% \begin{equation}
% \min_{\mbQ \in \reals^{N \times N}} \Vert \mbQ \mbX_i - \mbX_j \Vert_F \neq \min_{\mbQ \in \reals^{N \times N}} \Vert \mbX_i - \mbQ \mbX_j \Vert_F.
% \end{equation}

Finally, the Procrustes distance can be thought of as a constrained version of linear regression with an orthogonality constraint---i.e. a constraint such that the weight matrix must only encode a rotation and reflection of the firing rate space. Na\"ively dropping this restriction is problematic for constructing metric spaces. Indeed, for a generic $N \times N$ matrix of regression weights $\mbW$, the resulting distance score fails to even be symmetric.
% \begin{equation}
% \label{eq:naive-regression}
% \min_{\mbW_i} \Vert \mbW_i \mbX_i - \mbX_j \Vert_F \neq \min_{\mbW_j} \Vert \mbX_i - \mbW_j \mbX_j \Vert_F
% \end{equation}
This method of measuring similarity would therefore be sensitive to the order in which systems, such as animals or brain regions, are compared. In fact, simple forms of linear regression, even when symmetrized \parencite{muzellec2026reverse}, do not easily enable comparisons of neural systems within a metric space, as we show in simulation \hl{sup}.

% Of course, regressions and decoding analyses remain useful, complementary tools: precisely because they target a specific variable, they can reveal whether that information is present and readable by a downstream population—a question shape distances alone cannot answer. Indeed, we use decoding throughout this paper to validate and interpret the structure that shape distances reveal.

% In our experience, the orthogonality constraint inherent to Procrustes distance is often useful to regularize the alignment transformation.
% Nonetheless, this constraint can be relaxed by computing
% \begin{equation}
% \label{eq:whitened-procrustes-distance}
% d_{\text{Linear}}(\mbX_i, \mbX_j) = d_{\text{Proc}}(\mbZ_i \mbX_i, \mbZ_j \mbX_j)
% \end{equation}
% as the distance, where $\mbZ_i$ and $\mbZ_j$ are $N \times N$ \textit{whitening matrices}.
% This avoids the problem highlighted in \cref{eq:naive-regression} while also relaxing the orthogonality constraint inherent to Procrustes distance---i.e. we have that $d_{\text{Linear}}(\mbX_i, \mbX_j) = 0$ if any invertible linear transformation brings the two neural manifolds into alignment.
% The resulting distance is a proper metric and has connections to CCA, as discussed in \textcite{Williams2021}.

% Another, more stringent, metric we will utilize is the \textit{one-to-one matching distance}, which uses a permutation of the $N$ neurons to align the responses
% \begin{equation}
% \label{eq:one-to-one-matching}
% d_{\text{Perm}}(\mbX_i, \mbX_j) = \min_{\mbP \in \mbPi(N)} \Vert \mbP \mbX_i - \mbX_j \Vert_F
% \end{equation}
% where $\mbP$ is an $N \times N$ permutation matrix.
% Note that the set of permutation matrices, $\mbPi(N)$, is a subset of all orthogonal matrices, $\cO(N)$.
% Therefore, the Procrustes distance is a less stringent measure of similarity in the sense that $d_{\text{Proc}}(\mbX_i, \mbX_j) \leq d_{\text{Perm}}(\mbX_i, \mbX_j)$.
% \textcite{khosla2024soft} describe how to extend this distance to case where the two populations consist of unequal numbers of neurons.

% In this section, we surveyed several metrics that satisfy the desirable properties of \cref{eq:metric-space-equivalence,eq:metric-space-symmetry,eq:metric-space-triangle}.
% These belong to an even broader class of distances that are referred to as \textit{generalized shape metrics} \parencite{Williams2021}. 
% Additional mathematical details are provided in the \textit{\nameref{sec:methods}}.




% If we set $\phi(\mbX) = \mbX$ (i.e. no pre-processing) and allow $\mbQ$ to be any linear transformation, then equation~(\ref{eq:generalized-euclidean-shape-metric}) reduces to least-squares linear regression.
% This is not a metric; indeed it is not even symmetric. Somewhat remarkably, if we simply constrain $\mbQ$ to be an orthogonal matrix---i.e. implementing a rotation and reflection in $N$-dimensional space---the resulting distance is a bona fide metric.
% Formally, when we set $\cG$ equal to $\cO$ (the set of $N \times N$ orthogonal matrices) and $\phi(\mbX)$ performs mean-centering of each tuning curve, then equation (\ref{eq:generalized-euclidean-shape-metric}) becomes the \textit{Procrustes size-and-shape distance}~\parencite{Gower2004}.
% Intuitively, this distance treats $\mbX_i$ and $\mbX_j$ as equivalent if they have the same ``size-and-shape'' but potentially different orientations in neural firing rate space (\hl{figure}).

% Equation~(\ref{eq:generalized-euclidean-shape-metric}) encompasses many other possibilities beyond the Procrustes distance.
% For instance, we can recover a metric that is closely related to CCA when we choose $\phi(\mbX)$ to be a whitening transformation and optimize over orthogonal matrices (i.e., $\cG = \cO$, \hl{supplement}).
% This notion of distance is invariant to more general linear deformations, whereas Procrustes distance is only invariant to rotations in neural firing rate space.
% Another useful metric for neuroscience applications is to optimize over permutations of the neurons instead of orthogonal transformations.
% This too, produces a valid metric (\hl{supplement}), which intuitively treats two neural representations as equivalent only when one can find a one-to-one matching of the neurons.

% It is important to emphasize that Equation~(\ref{eq:generalized-euclidean-shape-metric}) does not specify a single metric, but rather encompasses a broad family of metrics that quantify different aspects of representational geometry and provide complementary scientific insights.
% We refer to these as \textit{generalized shape metrics}, since they generalize classical notions of shape distance \parencite{kendall2009shape}.
% Furthermore, in \hl{Supplementary Text} we show that measures based on RSA, CKA, and CCA are all related to this framework, and some can even be viewed as special cases of shape metrics after some modification.
% However, none of these existing approaches have been formulated in a way to satisfy and leverage all three criteria of a metric space. 



% Reducing each region to its first $10$ principal components, we computed the $K \times K$ matrix of Procrustes distances, $\mbD_{ij} = d_{\text{Proc}}(\mbX_i, \mbX_j)$, and embedded it by classical
% multidimensional scaling.
% Applying the same lumpy-versus-continuous test one level up, we compared the best mean silhouette over $k = 2 \dots 6$ in this embedding against $200$ Gaussians matched to it.
% The two cohorts give opposite answers, and in the direction opposite to their neuron-level answers.
% The first cohort, whose regions contain unambiguously categorical neurons, is indistinguishable from a single continuous cloud in region space ($p = 0.199$); what it contains instead is an ordering, with the
% first embedding dimension recovering position along the simulated gradient at $|\rho| = 0.96$ (\textbf{Fig. } \ref{fig:regionsim}c, d).
% The second cohort, in which not a single region is categorical, separates into three clean clusters far above the null ($p = 0.005$; \textbf{Fig. } \ref{fig:regionsim}g, h). 
  
%   This dissociation is not a matter of statistical power, and it is not repaired by putting regions in a common space by some simpler route.
%   Averaging each region's neurons into a single tuning curve, as is often done to summarise a region, and applying PCA across regions recovers neither structure: the gradient is lost ($|\rho| = 0.04$) and the three
%   types overlap (\textbf{Fig. } \ref{fig:regionsim}b, f). 
%   This is expected here by construction, since regions were built to share a common mean tuning curve, but it makes the general point concrete: an average over neurons retains the first moment of a region's tuning
%   distribution and discards its shape, which is precisely where the between-region structure lives.
%   Comparing populations as objects, with a distance that respects the metric axioms, is what makes the question answerable at all.


% As $\theta$ is varied, the firing rate traces out a smooth, 1D ring manifold in the high-dimensional firing rate space (\textbf{Fig. } \ref{fig:schematic}d).
% Figure~\ref{fig:synth}a shows tuning curves (\textit{left}) and PCA projections (\textit{right}) from three simulated subjects with $N=100$ neurons per subject.
% We simulated similar tuning curves from 37 additional subjects, resulting in a cohort size of $K=40$.

% In these simulations, we randomly varied two key variables across subjects.
% First, we varied the average tuning curve width.
% For example, network 1 has much sharper tuning curves than network 2 in \textit{Figure}~\ref{fig:synth}a.
% Second, we varied the concentration of tuning around $\theta=0$.
% In \textit{Figure}~\ref{fig:synth}a, network 3 has many more tuning curves peaked near the center of the horizontal axis.
% Non-uniform tiling of tuning curves is seen in neural circuits supporting navigation (e.g., over-representation of reward locations,~\cite{Hollup2001}) and sensory function (e.g., over-representation of cardinal edges in visual cortex,~\cite{Li2003}).

% We begin by computing a $K \times K$ matrix of pairwise Procrustes distances, with elements given by $\mbD_{ij} = d_{\text{Proc}}(\mbX_i, \mbX_j)$.
% % We will see that this matrix is a natural starting point to perform a wide variety of analyses.
% We first demonstrate how to use $\mbD$ to perform dimensionality reduction so that we can visualize the $K$ subjects as $K$ points in 2D space.
% We use multidimensional scaling (MDS; \cite{Borg2005-gk}), which is a well-established technique that takes in a matrix of pairwise distances, $\mbD$, and returns a list of vectors, $\mbu_1, \dots, \mbu_K$, such that $\Vert \mbu_i - \mbu_j \Vert_2 \approx \mbD_{ij}$ for all pairs of subjects.
% Some amount of mismatch between the distances is unavoidably introduced by the curvature of the underlying shape space~\parencite{Robinson2006-vn,kendall2009shape}.
% However, the distortion is typically small, especially if the embedded vectors, $\mbu_1, \dots, \mbu_K$, are allowed to occupy a high-dimensional space.

% In our simulated example, a 20-dimensional MDS analysis produced a median distortion of \textapprox{}3.4\% across all pairwise distances (see \textit{\nameref{sec:methods}}).
% A 2D PCA projection of this embedding, capturing \textapprox{}85\% of the variance, is shown in \textit{Figure}~\ref{fig:synth}b-c.
% From these visualizations, we see the shape metric successfully reveals two forms of across-subject variability that we built into the simulation: subjects are ordered by tuning curve width along one dimension (\textbf{Fig. } \ref{fig:synth}b, top to bottom) and ordered by their concentration around $\theta=0$ along another (\textbf{Fig. } \ref{fig:synth}c, left to right).
% Thus, variability in tuning curve parameters across subjects is transparently captured within the Procrustes shape metric space. Here, we explore variability across subjects using Procrustes distance, which is invariant to rotations of the neural firing-rate space. This invariance has an important consequence: two populations can share the same population-level geometry while differing completely in the tuning of individual neurons, because a rotation rearranges single-neuron tuning without altering the shape of the manifold (\cite{Kriegeskorte2021}). We illustrate this by generating a randomly rotated copy of each simulated network: single-neuron tuning is scrambled (\textbf{Fig. }~\ref{fig:synth}d), yet each copy is indistinguishable from its original in Procrustes space and lies directly on top of it (\textbf{Fig. }~2f). This is precisely the level of description we seek: population geometry abstracted from the identity of individual neurons.

% In our simulation, the kNN performance is improved by increasing the total number of networks, $K$, as well as the number of recorded neurons per network, $N$.
% Intuitively, increasing $K$ improves the model by creating a denser sampling of the Procrustes metric space, while increasing $N$ improves our estimate of each subject's position in the Procrustes metric space.

% A simple approach to accomplish this is to fit use an Euclidean embedding of shape space (found by MDS) to fit a regression model.
% This works very well on this simulated dataset: we can predict a subject's performance (i.e. linear fisher information) with $R^2 = $\hl{xxx} (evaluated on heldout data, leave-one-out basis).
% We emphasize that this behavioral prediction is made \textit{solely using the held out subject's neural data}---a much more difficult task than jointly training a decoder on behavioral and neural data.

% MDS will always introduce some level of distortion into the underlying metric space~\parencite{Robinson2006-vn}.
% Therefore, it is arguably more elegant to perform regression directly on the shape space, which can be done using a nonparametric nearest neighbors model~\parencite{Cover1967} or various geodesic regression methods~\parencite{something}.
% We review and demonstrate these more advanced approaches in \hl{supplement} and show that they produce qualitatively similar insights.

% Above, we explored variability across subjects using a rotation-invariant shape metric (Procrustes distance).
% However, this represents only one choice from the broader family of generalized shape metrics.
% To demonstrate this, we generate a randomly rotated copy of each network.
% This results in a neural representations with the same geometry but completely rearranged tuning (\textbf{Fig. } \ref{fig:synth}e).
% By construction, each rotated copy is indistinguishable from its original counterpart in terms of Procrustes distance.
% Thus, when we visualize all of the simulated subjects within Procrustes shape space, we see that each neural system and its rotated copy lie on top of each other (\textbf{Fig. } \ref{fig:synth}f).

% The Procrustes distance does not necessarily detect differences in the structure of the tuning curve, since many shapes of the tuning curve can give rise to the same geometry at the population level~\parencite{Kriegeskorte2021}.
% However, other shape metrics may be used to quantify and detect these differences.
% To show this, we use the metric in \cref{eq:one-to-one-matching} that is invariant to permutations of the neuron labels but sensitive to rotations (see \textit{\nameref{sec:methods}}).
% When we repeat the procedure above with this new distance measure, we see that the rotated copies are positioned far away from their original counterparts (\textbf{Fig. } \ref{fig:synth}g).
% These two clusters (original vs. rotated networks) are also readily revealed by hierarchical clustering applied to the $K \times K$ distance matrix.

% The purpose of this simulation is to show that different shape metrics are suited to different scientific questions.
% Thus, depending on the broader context, the insensitivity of Procrustes distance to rotations can be construed either as a weakness or a strength.
% To show this, we revisit the problem of using the $K \times K$ shape distance matrix to infer the behavioral performance (Fisher information) of each simulated subject.
% Since the Fisher information is invariant to rotations, each network has the same information score as its randomly rotated counterpart.
% Thus, when we use a $k$-nearest neighbor model to predict the information score from the permutation-invariant distance matrix, we observe substantially worse performance ($R^2=0.37$; \textbf{Fig. } \ref{fig:synth}h blue markers) than when the same model is fit on Procrustes distances ($R^2=0.83$; \textbf{Fig. } \ref{fig:synth}h black markers).

%\subsection{Relating behavioral differences to position in neural shape space}

% Now we show how to leverage these distance relationships to make predictions about individual subjects.
% Suppose that we are interested in each subject's ability to perform a sensory discrimination task between two conditions, ${\theta = \pm \epsilon}$ where ${\epsilon \approx 0}$.
% A popular way to quantify the ability of a neural representation to support perceptual discrimination is \textit{Fisher Information} (FI), a quantity that is related to manifold geometry and structure of noise \parencite[see e.g.,][]{Kriegeskorte2021}.
% We computed FI per neuron as a proxy for each subject's acuity, or theoretical performance on a discrimination task (see \textit{Methods}), and then asked whether we could predict this acuity score for each subject using only their position within the Procrustes metric space.
% We found that a $k$-nearest neighbor (kNN) regression model~\parencite{Cover1967} that predicts each subject's acuity by averaging over its $k$ nearest neighbors in Procrustes shape space accurately predicted acuity scores in this simulation ($R^2 = 0.95$, $k=3$).

